% SPDX-License-Identifier: CC-BY-SA-4.0
% Author: Matthieu Perrin

\begingroup

\begin{frame}[fragile]{Le consensus}

  \begin{block}{Définition}
    \begin{lstlisting}
      interface Consensus<T> {
        public T propose(T val) ;
      }
    \end{lstlisting}
    Chaque thread \structure{propose} une valeur et peut en \structure{décider} une.
    \begin{description}[Wait-freedom : ]
    \item[Accord : ] Au plus une valeur décidée
    \item[Validité : ] Toute valeur décidée a été proposée
    \item[Wait-freedom : ] Tout thread correct décide
    \end{description}
  \end{block}

  \pause

  \begin{exampleblock}{Feuille de route}
    \begin{description}[Thm1 :]
    \item[Thm1 :] Si on pouvait implémenter une séquence, \\ \hspace{5mm} alors on pourrait résoudre le consensus {\color{exampleColor}entre deux threads}
    \item[Thm2 :] Or, on ne peut pas résoudre le consensus entre deux threads
    \item[Conc :] Donc on ne peut pas implémenter une séquence
    \end{description}
  \end{exampleblock}

\end{frame}

\endgroup
\endinput
